\documentclass[11pt,a4paper]{article}
\usepackage[utf8]{inputenc}
\usepackage[french]{babel}
\usepackage[T1]{fontenc}
\usepackage{geometry}
\usepackage{multicol}
% Pas de package enumitem
\usepackage{titlesec}
\usepackage{listings}
\usepackage{xcolor}

% Configuration de la page (Marges plus confortables)
\geometry{top=1.5cm, bottom=1.5cm, left=1.5cm, right=1.5cm}
\setlength{\parindent}{0pt}
\setlength{\parskip}{0.5em} % Espacement entre les paragraphes

% Réglage manuel des listes (remplacement enumitem)
\setlength{\leftmargini}{1.5em} 

% Couleurs du code
\definecolor{codebg}{rgb}{0.96,0.96,0.96}
\definecolor{keyword}{rgb}{0.1,0.1,0.7}
\definecolor{string}{rgb}{0.7,0.1,0.1}
\definecolor{comment}{rgb}{0.3,0.6,0.3}

% Configuration de l'affichage du code
\lstset{
    language=HTML,
    basicstyle=\footnotesize\ttfamily, % Taille plus lisible
    backgroundcolor=\color{codebg},
    keywordstyle=\color{keyword}\bfseries,
    commentstyle=\color{comment},
    stringstyle=\color{string},
    breaklines=true,
    frame=leftline,
    framesep=5pt,
    showstringspaces=false,
    tabsize=2,
    extendedchars=true,
    morekeywords={let, const, var, function, class, constructor, this, return, if, else, for, while, switch, case, break, new, extends, super, try, catch, document, window, console}
}

% Styles des titres
\titleformat{\section}{\large\bfseries\uppercase}{}{0em}{}[\hrule]
\titleformat{\subsection}{\normalsize\bfseries\color{blue}}{}{0em}{}
\titlespacing{\section}{0pt}{10pt}{5pt}
\titlespacing{\subsection}{0pt}{10pt}{2pt}

\begin{document}

% Titre du document
\begin{center}
    {\Huge \textbf{Fiche Révision : HTML5 / CSS3 / JS}}
\end{center}

\begin{multicols*}{2}

% ==========================================
% PARTIE 1 : HTML5
% ==========================================
\section{I. HTML5 : Structure \& Sémantique}

\subsection{Squelette standard}
\begin{lstlisting}
<!DOCTYPE html>
<html lang="fr">
<head>
    <meta charset="utf-8">
    <meta name="viewport" content="width=device-width, initial-scale=1.0">
    <title>Titre de la page</title>
    <link rel="stylesheet" href="style.css">
    <script src="app.js" defer></script>
</head>
<body>
    </body>
</html>
\end{lstlisting}

\subsection{Principales balises sémantiques}
Ces balises structurent le document pour le référencement et l'accessibilité.
\begin{itemize}
    \item \texttt{<header>} : En-tête de page ou de section.
    \item \texttt{<nav>} : Liens de navigation principaux.
    \item \texttt{<main>} : Contenu principal unique de la page.
    \item \texttt{<section>} : Regroupement thématique de contenu.
    \item \texttt{<article>} : Contenu autonome (ex: post de blog).
    \item \texttt{<aside>} : Contenu indirectement lié (barre latérale).
    \item \texttt{<footer>} : Pied de page (copyright, contacts).
\end{itemize}

\subsection{Formulaires}
L'attribut \texttt{name} est indispensable pour l'envoi des données (PHP/Back-end).
\begin{lstlisting}
<form action="traitement.php" method="POST">
    <label for="nom">Nom :</label>
    <input type="text" id="nom" name="nom" placeholder="Votre nom" required>

    <input type="email" name="mail">
    <input type="password" name="pass">

    <input type="radio" name="genre" value="F"> Femme
    <input type="radio" name="genre" value="H"> Homme
    
    <select name="pays">
        <option value="fr">France</option>
        <option value="be">Belgique</option>
    </select>

    <textarea name="message" rows="5"></textarea>

    <input type="submit" value="Envoyer">
</form>
\end{lstlisting}

\subsection{Tableaux}
\begin{lstlisting}
<table>
    <thead>
        <tr> <th>Titre 1</th> <th>Titre 2</th> </tr>
    </thead>
    <tbody>
        <tr> <td>Data A</td> <td>Data B</td> </tr>
    </tbody>
</table>
\end{lstlisting}

\subsection{Médias}
\begin{lstlisting}
<img src="img.jpg" alt="Description">
<a href="page.html" target="_blank">Lien</a>
\end{lstlisting}

% ==========================================
% PARTIE 2 : CSS3
% ==========================================
\section{II. CSS3 : Mise en forme}

\subsection{Sélecteurs courants}
\begin{itemize}
    \item \texttt{p} : Cible toutes les balises \texttt{<p>}.
    \item \texttt{.maclasse} : Cible les éléments avec \texttt{class="maclasse"}.
    \item \texttt{\#monid} : Cible l'élément unique avec \texttt{id="monid"}.
    \item \texttt{div p} : Les \texttt{p} descendants de \texttt{div}.
    \item \texttt{div > p} : Les \texttt{p} enfants directs de \texttt{div}.
    \item \texttt{input[type="text"]} : Cible par attribut.
\end{itemize}

\subsection{Modèle de boîte (Box Model)}
\begin{lstlisting}
.box {
    width: 300px;       /* Largeur contenu */
    padding: 20px;      /* Marge interne */
    border: 1px solid;  /* Bordure */
    margin: 10px;       /* Marge externe */
    box-sizing: border-box; /* Padding inclus dans width */
}
\end{lstlisting}

\subsection{Flexbox (Mise en page 1D)}
S'applique sur le conteneur parent pour aligner les enfants.
\begin{lstlisting}
.container {
    display: flex;
    /* Direction : row (defaut) | column */
    flex-direction: row; 
    
    /* Axe Principal (horizontal si row) */
    justify-content: space-between; /* center, flex-start... */
    
    /* Axe Secondaire (vertical si row) */
    align-items: center; /* stretch, flex-start... */
    
    /* Retour a la ligne si manque de place */
    flex-wrap: wrap; 
}
\end{lstlisting}

\subsection{CSS Grid (Mise en page 2D)}
\begin{lstlisting}
.grille {
    display: grid;
    /* 3 colonnes : 1fr, 2fr et 1fr */
    grid-template-columns: 1fr 2fr 1fr; 
    gap: 20px; /* Espacement */
}
.item-a {
    grid-column: 1 / 3; /* Occupe col 1 a 3 */
}
\end{lstlisting}

% ==========================================
% PARTIE 3 : JAVASCRIPT
% ==========================================
\section{III. JavaScript : Algorithmique}

\subsection{Déclaration de variables}
\begin{lstlisting}
let score = 0;        // Variable modifiable (portee bloc)
const PI = 3.14;      // Constante (non reassignable)
var old = "vieux";    // Ancienne syntaxe (portee fonction)
\end{lstlisting}

\subsection{Structures de contrôle}
\begin{lstlisting}
// Condition
if (age >= 18) {
    console.log("Majeur");
} else if (age > 12) {
    console.log("Ado");
} else {
    console.log("Enfant");
}

// Switch
switch(action) {
    case 'START': demarrer(); break;
    default: arret();
}

// Ternaire
let statut = (age >= 18) ? "Adulte" : "Mineur";
\end{lstlisting}

\subsection{Boucles}
\begin{lstlisting}
// For classique (nombre d'iterations connu)
for (let i = 0; i < 5; i++) {
    console.log(i);
}

// While (tant que condition vraie)
while (compteur < 10) {
    compteur++;
}

// Parcours de tableau
const fruits = ['Pomme', 'Poire', 'Banane'];
fruits.forEach(fruit => {
    console.log(fruit);
});
\end{lstlisting}

\subsection{Fonctions}
\begin{lstlisting}
// Fonction nommee classique
function somme(a, b) {
    return a + b;
}

// Fonction flechee (Arrow Function - ES6)
const carre = (x) => {
    return x * x;
};
// Syntax courte si une seule ligne :
const double = x => x * 2;
\end{lstlisting}

% ==========================================
% PARTIE 4 : DOM & EVENTS
% ==========================================
\section{IV. JavaScript : DOM \& Events}

\subsection{Sélection d'éléments (DOM)}
\begin{lstlisting}
// Selection unique (premier trouve)
const titre = document.getElementById('main-title');
const btn = document.querySelector('.btn-primary');

// Selection multiple (NodeList)
const liens = document.querySelectorAll('nav a');
\end{lstlisting}

\subsection{Manipulation du contenu et du style}
\begin{lstlisting}
const boite = document.querySelector('#box');

// Lecture/Ecriture texte HTML
boite.innerHTML = "<strong>Nouveau</strong> contenu";
boite.textContent = "Texte brut uniquement";

// Classes CSS
boite.classList.add('active');
boite.classList.remove('hidden');
boite.classList.toggle('dark-mode');

// Style direct
boite.style.backgroundColor = "red";
boite.style.display = "none";

// Attributs HTML
boite.setAttribute('title', 'Info-bulle');
console.log(boite.getAttribute('href'));
\end{lstlisting}

\subsection{Gestion des événements (Events)}
Méthode recommandée : \texttt{addEventListener}.

\textbf{Principaux événements :}
\begin{itemize}
    \item \texttt{click}, \texttt{dblclick} (Souris)
    \item \texttt{mouseover}, \texttt{mouseout} (Survol)
    \item \texttt{input}, \texttt{change}, \texttt{submit} (Formulaires)
    \item \texttt{keydown}, \texttt{keyup} (Clavier)
    \item \texttt{DOMContentLoaded}, \texttt{load} (Chargement page)
\end{itemize}

\begin{lstlisting}
const bouton = document.querySelector('#monBouton');

bouton.addEventListener('click', function(event) {
    // Empeche le comportement par defaut (ex: envoi formulaire)
    event.preventDefault(); 
    
    alert("Bouton clique !");
    console.log(event.target); // L'element clique
});
\end{lstlisting}

\subsection{Création d'éléments dynamiques}
\begin{lstlisting}
// 1. Creation
const newDiv = document.createElement('div');
newDiv.textContent = "Je suis nouveau";

// 2. Insertion
const parent = document.querySelector('body');
parent.appendChild(newDiv); 

// Suppression
// parent.removeChild(newDiv);
\end{lstlisting}

% ==========================================
% PARTIE 5 : JS AVANCÉ (Objets)
% ==========================================
\section{V. JavaScript : Objets \& Tableaux}

\subsection{Objet Littéral}
\begin{lstlisting}
const user = {
    nom: "Dupont",
    prenom: "Jean",
    age: 25,
    direBonjour: function() {
        return "Salut " + this.prenom;
    }
};
console.log(user.nom); // Acces propriete
\end{lstlisting}

\subsection{Tableaux et méthodes utiles}
\begin{lstlisting}
let tab = [10, 20, 30];

tab.push(40);     // Ajoute a la fin
tab.pop();        // Supprime dernier
tab.length;       // Taille du tableau (3)
tab.splice(0, 1); // Supprime 1 element a l'index 0
\end{lstlisting}

\subsection{Classes (POO ES6)}
\begin{lstlisting}
class Vehicule {
    constructor(marque, couleur) {
        this.marque = marque;
        this.couleur = couleur;
    }
    
    demarrer() {
        console.log("Vrooom");
    }
}

const maVoiture = new Vehicule("Peugeot", "Rouge");
maVoiture.demarrer();
\end{lstlisting}

\end{multicols*}
\end{document}